\section{Практическая часть}
В листинге~\ref{lst:source} представлен исходный код реализованной программы.
\begin{listing}[Исходный код программы]{lst:source}{Matlab}
clear; clc; close all;

X = [7.76,6.34,5.11,7.62,8.84,4.68,8.65,6.90,8.79,6.61, ...
    6.62,7.13,6.75,7.28,7.74,7.08,5.57,8.20,7.78,7.92, ...
    6.00,4.88,6.75,6.56,7.48,8.51,9.06,6.94,6.93,7.79, ...
    5.71,5.93,6.81,5.76,5.88,7.05,7.22,6.67,5.59,6.57, ...
    7.28,6.22,6.31,5.51,6.69,7.12,7.40,6.86,7.28,6.82, ...
    7.08,7.52,6.81,7.55,4.89,5.48,7.74,5.10,8.17,7.67, ...
    7.07,5.80,6.10,7.15,7.88,9.06,6.85,4.88,6.74,8.76, ...
    8.53,6.72,7.21,7.42,8.29,8.56,9.25,6.63,7.49,6.67, ...
    6.79,5.19,8.20,7.97,8.64,7.36,6.72,5.90,5.53,6.44, ...
    7.35,5.18,8.25,5.68,6.29,6.69,6.08,7.42,7.10,7.14, ...
    7.10,6.60,6.35,5.99,6.17,9.05,6.01,7.77,6.27,5.81, ...
    7.80,9.89,4.39,6.83,6.53,8.15,6.68,6.87,6.31,6.83]

X = X(:); 
n = length(X);

M_min = min(X);
M_max = max(X);
R = M_max - M_min;

mu_hat = mean(X);
S2 = var(X);
S = sqrt(S2);

m = floor(log2(n)) + 2; 
edges = linspace(M_min, M_max, m + 1);

fprintf('--- Результаты расчетов ---\n');
fprintf('Объем выборки        = %d\n', n);
fprintf('Минимум              = %.4f\n', M_min);
fprintf('Максимум             = %.4f\n', M_max);
fprintf('Размах               = %.4f\n', R);
fprintf('Выборочное среднее   = %.4f\n', mu_hat);
fprintf('Исправленная выборочная дисперсия = %.4f\n', S2);
fprintf('Количество интервалов m = %d\n\n', m);

x_grid = linspace(M_min - 0.1*R, M_max + 0.1*R, 1000);

figure(1);
hold on; grid on;

histogram(X, 'BinEdges', edges, 'Normalization', 'pdf', 'FaceAlpha', 0.6, 'FaceColor', '#0072BD');

y_pdf = (1 / (S * sqrt(2 * pi))) * exp(-((x_grid - mu_hat).^2) / (2 * S^2));
plot(x_grid, y_pdf, 'LineWidth', 2, 'Color', 'r');

title('Гистограмма и плотность нормального распределения');
xlabel('Значения X');
ylabel('Плотность вероятности f(x)');
legend('Эмпирическая плотность (Гистограмма)', 'Плотность N(\mu, S^2)');

figure(2);
hold on; grid on;

X_sorted = sort(X);
x_emp =[M_min - 0.1*R; X_sorted];
F_emp =[0; (1:n)' / n];

stairs(x_emp, F_emp, 'LineWidth', 2, 'Color', 'b');

y_cdf = 0.5 * (1 + erf((x_grid - mu_hat) / (S * sqrt(2))));
plot(x_grid, y_cdf, 'LineWidth', 2, 'Color', 'r', 'LineStyle', '--');

title('Эмпирическая и теоретическая функции распределения');
xlabel('Значения X');
ylabel('Функция распределения F(x)');
legend('Эмпирическая ФР', 'Функция распределения N(\mu, S^2)', 'Location', 'NorthWest');
\end{listing}

В листинге~\ref{lst:calc} приведены результаты расчетов указанных величин.
\begin{listing}[Результаты расчетов указанных величин]{lst:calc}{Matlab}
Объем выборки           = 120
Минимум                 = 4.3900
Максимум                = 9.8900
Размах                  = 5.5000
Выборочное среднее      = 6.9445
Исправленная выборочная дисперсия = 1.1720
Количество интервалов m = 8
\end{listing}

На рисунке~\ref{img:f1.pdf} представлены гистограмма и график функции плотности 
распределения вероятностей нормальной случайной величины с математическим 
ожиданием \( \hat{\mu} \) и дисперсией \( S^2 \).

\jimg{f1.pdf}{Гистограмма и график функции плотности распределения вероятностей
нормальной случайной величины}

На рисунке~\ref{img:f2.pdf} представлены график эмпирической функции 
распределения и график функции распределения нормальной случайной величины 
с математическим ожиданием \( \hat{\mu} \) и дисперсией \( S^2 \).

\jimg{f2.pdf}{Графики эмпирической функции распределения и функции распределения
нормальной случайной величины}
